\documentclass[12pt]{article}
\usepackage{amsmath}
\usepackage{bm}
\usepackage{graphicx}
\usepackage{geometry}
\usepackage{indentfirst}
\usepackage{amsfonts}
\geometry{legalpaper, portrait, margin=0.5in}
\usepackage{color}   %May be necessary if you want to color links
\usepackage{hyperref}
\hypersetup{
    colorlinks=true, %set true if you want colored links
    linktoc=all,     %set to all if you want both sections and subsections linked
    linkcolor=black,  %choose some color if you want links to stand out
}
\usepackage[siunitx]{circuitikz}
\usetikzlibrary{patterns}
\usetikzlibrary{decorations.markings}
\graphicspath{ {./images/} }
\usepackage{enumitem}
\usepackage{outlines}
\begin{document}
\newcommand*\dif{\mathop{}\!\mathrm{d}}
\renewcommand{\mod}{\,\operatorname{mod}\,}
\setlength{\parindent}{0pt}

\setlist[enumerate]{noitemsep, topsep=0pt, itemindent=\parindent}
\setlist[itemize]{noitemsep, topsep=0pt, itemindent=\parindent}

\newenvironment{subitemize}
{ \begin{itemize}[leftmargin=\parindent] }
{ \end{itemize} }

\newenvironment{subenumerate}
{ \begin{enumerate}[leftmargin=\parindent] }
{ \end{enumerate} }

\newenvironment{nopagebr}
  {\par\nobreak\vfil\penalty0\vfilneg
   \vtop\bgroup}
  {\par\xdef\tpd{\the\prevdepth}\egroup
   \prevdepth=\tpd}

\newcommand{\bbr}{\mathbb{R}}

\pagebreak
\begin{outline}

\section{Problem}

We've done a forward pass in a neural network $f_\theta$. But backprop for parameter update needs $\nabla_\theta J$, and that's a lot of derivation via chain rule. We want to automate this --- after all, it's just differentiation. \\

However, the limit definition of the derivative $\displaystyle \lim_{h \rightarrow 0} \frac{f(x+h)-f(x)}{h}$ for computing $\nabla_\theta J$ has some issues: it's numerically unstable (floating point errors for small $h$) and slow (each parameter $\theta_i$ requires its own full forward pass through the network). We need something faster and more reliable.

\section{Solution}

\underline{\textbf{Big-picture idea}}: we'll hard-code a set of primitive functions along with its derivatives with respect to each of its arguments: a primitive function $f(a_0,...,a_{n-1})$ hard-codes $n$ derivative functions (one for each positional argument) $d_{f,i} = \displaystyle \frac{\partial f}{\partial a_i}$ for each $i$. Then, we'll represent bigger functions (like $J(\theta)$, which involves a full forward pass and then a cost function on the network's final output) as simply compositions of these primitive functions; the gradient $\nabla_\theta J$ is computable via chain rule across these smaller primitives. \\

Construct a computational graph that's centered around primitive functions: each node represents a function call.
\1 Each node $u$ owns:
\2 A \textbf{primitive function} $f_u$ (which primitive function it is: addition, multiplication, relu --- these must be hard-coded)
\2 An \textbf{arg list} $\text{adj}(u)$ (the nodes whose results are used as arguments for this node's fn) --- this is the node's adjacency list, store edges in order of the positional args.
\2 A cached \textbf{result} $\hat{u}$ (the value as a result of its function taking in its args)
\2 A cached \textbf{gradient} $\text{grad}(u)$ (which will be used in backprop; if this node represents $\hat u = f_u(\hat{\text{adj}(u)})$, then the gradient value holds $\displaystyle \frac{\partial \text{root}}{\partial \hat u}$)
\1 This will form a DAG where the cost function call $J$ roots the DAG, and the parameters $\theta_i$ are the leaves of the DAG.
\\\0

The graph and its nodes are all constructed as computation occurs. Then, we can start from any node $s$ and compute $\displaystyle \frac{\partial \hat s}{\partial \hat u}$ for all $u$ reachable from $s$. We must traverse the graph in topological order when computing gradients.
\1 In our neural network case, we'd want to start from node $s$ where $J=\hat s$ and end at all nodes $t_i$ where $\theta_i = \hat t_i$, in order to compute $\nabla_\theta J$. Obviously, $\displaystyle \frac{\partial J}{\partial J} = 1$ to start, and initialize all grad values to zero.
\1 Then, suppose we're sitting at node $v$. For each edge that exists in $\text{adj}(v)$, let the child be $u$, and suppose it's specifically passed in as the $i$th positional argument which is denoted as $a_i = \hat{u}$. $a_i$ only appears as the $i$th argument, even if $u$ appears among multiple positional args; $\displaystyle \frac{\partial J}{\partial \hat{u}} += \sum_{a_i} \frac{\partial J}{\partial a_i}$ for a single parent by multivariate chain rule, and we do this for all parents of $u$ for the full $\displaystyle \frac{\partial J}{\partial \hat u}$.
\2 We need to find $\displaystyle \frac{\partial J}{\partial a_i} = \frac{\partial J}{\partial \hat v} \frac{\partial \hat v}{\partial a_i}$. Equivalently: $\displaystyle \frac{\partial J}{\partial a_i} = \text{grad}(v) \frac{\partial \hat v}{\partial a_i}$. (notice: needing $\displaystyle \frac{\partial J}{\partial \hat v}$ to be finalized is why we need topo traversal --- $v$ can't explore children before $\text{grad}(v)$ is finalized.) But what's $\displaystyle \frac{\partial \hat v}{\partial a_i}$?
\2 Remember that $\hat v = f_v(\hat{\text{adj}(v)})$, so differentiating both sides w.r.t the positional arg, we get
\[ \frac{\partial \hat v}{\partial a_i} = \frac{\partial}{\partial a_i} f_v(\hat{\text{adj}(v)}) = d_{f_v, i}(\hat{\text{adj}(v)})\]
\2 So then in the end, we do $\displaystyle \boxed{\text{grad}(u) += \text{grad}(v) \cdot d_{f_v, i}(\hat{\text{adj}(v)})}$.
\1 We traverse until we've finished at the leaves. Then, we pick out the gradients of the nodes we care about (aka the nodes $t_i$ such that $\hat t_i = \theta_i$, which gives us $\displaystyle \frac{\partial J}{\partial \theta_i}$).
\0

\section{Specializing for matmuls}

The core operation in a neural net --- matmuls --- generate a gigantic autograd graph. Our primitives are just too simple. We need to elevate them to represent larger primitive operations. We'll have to make some optimizations; if we naively try $C = \text{matmul}(A,B) = AB$, then $d_{\text{matmul},i}$ gives us a giant 4D Jacobian, where $\displaystyle \frac{\partial C}{\partial A}$ for example would have a matrix of $\displaystyle \frac{\partial C_{ij}}{\partial A_{i' j'}}$ for all $i',j'$ given a fixed $i,j$ --- a matrix of partials for \textbf{each} entry in $C$. \\

So we'll just compute the chain rule directly, to avoid constructing such a Jacobian. If you do the algebra and simplify, we can define a new function $g$ that directly gives us the gradient wrt the positional arg we want (denoting gradient matrix $\displaystyle G = \frac{\partial J}{\partial C}$):
\1 $\displaystyle \frac{\partial J}{\partial A} = g_{\text{matmul},0}(G,A,B) = GB^t$
\1 $\displaystyle \frac{\partial J}{\partial B} = g_{\text{matmul},1}(G,A,B) = A^t G$
\\\0

So in practice, if you're doing tensor autograd, replace $d_{f_v,i}(\text{pos-args})$ with $g_{f_v,i}(\text{grad}(v), \text{pos-args})$ --- if visiting child $u$ from parent $v$, do $\boxed{\text{grad}(u) += g_{f_v,i}\left(\text{grad}(v), \hat{\text{adj}(v)}\right)}$.

\end{outline}
\end{document}
